\documentclass{article}
\usepackage[T1]{fontenc}
\usepackage{lmodern}
\usepackage{graphicx}
\usepackage{amsmath,amssymb}
\usepackage[a4paper,margin=2cm]{geometry}
\usepackage[absolute,overlay]{textpos}
\setlength{\TPHorizModule}{1mm}
\setlength{\TPVertModule}{1mm}
\usepackage[indonesian]{babel}
\usepackage{hyperref}
\setlength{\emergencystretch}{2em}
% unit: Halaman 37

\begin{document}

% === Halaman 37 (llm) ===

Secara matematis, mode $CFB$ $s$-bit dapat dinyatakan sebagai:

\begin{align*}
\text{Enkripsi:}\\ C_j &= P_j \oplus \text{MSB}_s(E_K(X_j)) \\
X_{j+1} &= \text{LSB}_{b-s}(X_j) \parallel C_i
\end{align*}
\begin{align*}
\text{Dekripsi:}\\ P_j &= C_i \oplus \text{MSB}_s(E_K(X_j)) \\
X_{j+1} &= \text{LSB}_{b-s}(X_j) \parallel C_i
\end{align*}

yang dalam hal ini,
\begin{itemize}
    \item $j = 1, 2, 3, ..., m$
    \item $X_j = \text{antrian bergeser, dengan } X_1 = IV$
    \item $E = \text{fungsi enkripsi}$
    \item $K = \text{kunci}$
    \item $b = \text{panjang blok enkripsi}$
    \item $s = \text{panjang unit enkripsi}$
    \item $\parallel = \text{operator penyambungan (concatenation)}$
    \item $\text{MSB} = \text{Most Significant Byte}$
\end{itemize}

\end{document}
